%!TEX root =  tfg.tex
\chapter{Metodología de desarrollo}
 
\begin{quotation} [Escritor] {Kevin Kelly (1952-- )}
Una de las funciones de una organización, de cualquier organismo, es anticipar el futuro, de modo que esas relaciones pueden persistir en el tiempo.
\end{quotation}

\begin{abstract}
Para realizar el proyecto ha sido necesario seguir una serie de procedimientos para el correcto desarrollo del mismo.
\end{abstract}

\section{Estructura organizacional del proyecto}

El proyecto en su totalidad tal y como lo hemos mencionado anteriormente se desarrollará en una pareja, pero cada esta estará partida en dos grandes módulos principales siendo estos el módulo de permutas (centrado en la búsquda y generación de intercambios de grupos) y el módulo del aseguramiento del servicio.

En este caso Francisco José Anillo se encargaría del módulo de permutas mientras que su compañero Daniel Dorante Sánchez del módulo de aseguramiento del servicio.

\section{Metodología de desarrollo}

En cuanto a las metodologías que se han seguido para el desarrollo del proyecto, se han definido varias metodologías para el desarrollo del proyecto y la gestión de la programación del proyecto.

También se han de definir ciertas políticas de gestión del código fuente, entre los que se encuentran, 

\subsection{Estrategia de ramas}

Comenzando por la estrategia de ramas se define de la siguiente manera:

La estrategia para la creación, uso, nomenclatura y cualquier otro aspecto que se encontrará basada en GitFlow, de forma que se constará con una rama de desarrollo (nombre particular "develop"), A partir de estas se crearan el resto dependiendo de cada funcionalidad que se vaya a desarrollar.

La nomenclatura de todas las ramas a crear mantendrá la siguiente estructura:

[Tipo] - [Numero]

\begin{itemize}
    \item Tipo: es el tipo de rama que se va a crear puede ser \textbf{feacture} o \textbf{hotfix} o \textbf{test}
    \item Numero: es el número de la issue que se implementa.
\end{itemize}

\subsection{Política de commits}

Seguiremos una política de commits en el que una rama atenderá a una issue. De esta forma se creará una única tarea por cada issue en el tablero de desarrollo.

La estructura de los commits responderá a la siguiente fórmula:

    [TituloFuncionalidad]

[DescripciónModificación]

Donde:
\begin{itemize}
    \item TituloFuncionalidad: es un título que describe la funcionalidad que se desarrolla.
    \item DescripciónModificación: es una descripción resumida de lo que se ha integrado por medio de los commits.
\end{itemize}
